\documentclass[10pt,a4paper]{amsart}
\usepackage[margin=19mm]{geometry}
\usepackage{amsmath,amssymb,mathtools}
\usepackage[scr=rsfs,cal=euler]{mathalfa}
\usepackage{xfrac}
\usepackage[unicode,hidelinks,bookmarks=false]{hyperref}
\newcommand{\ie}{i.e.\ }
\newcommand{\cf}{cf.\ }
\newcommand{\resp}{respectively\ }
\newcommand{\Subsection}{\subsection}
\providecommand{\nobreakdash}{\nobreak\hbox{-}}
\setlength{\emergencystretch}{2em}
\multlinegap0pt
\usepackage{amssymb}
\usepackage{xcolor}
\usepackage{enumerate}
\newtheorem{theorem}{Theorem}
\newtheorem{lemma}{Lemma}
\title{On the number of modular pairs\\ in finite dimensional Lie algebras on finite fields}
\author{Seid Kassaw Muhie, Daniele Ettore Otera, Francesco G. Russo}
\date{Source: arXiv:2609.19086v1 (CC BY 4.0)}
\begin{document}
\maketitle

\noindent\emph{Source and licence.} On the number of modular pairs\\ in finite dimensional Lie algebras on finite fields. Version arXiv:2609.19086v1 (CC BY 4.0), distributed by arXiv under the Creative Commons Attribution 4.0 International licence, http://creativecommons.org/licenses/by/4.0/, as recorded in the arXiv licence metadata for this exact version and shown on the abstract page https://arxiv.org/abs/2609.19086v1. Full licence notice: \texttt{licenses/arxiv-2609.19086-CC-BY-4.0.txt}.\par\smallskip

\noindent\textit{Editorial scope.} Counted unit: the article's theorem sdforsl2 (source0.tex lines 131--133). The article's complete original proof of it is delivered with it (source0.tex lines 529--579, one \texttt{proof} environment of 51 lines, which the article places away from the statement). No mathematics is written, repaired or re-derived: the statement and the proof are the article's own lines, in the article's own notation, and no retained line is substituted or re-flowed.  Retained uncounted as the source-local context the proof needs: lines 124--127; lines 280--282; lines 304--308.  Nothing was omitted from any retained span, so the package carries no omission record.  No retained line contains a citation, so the wrapper carries no bibliography and no bibliography entry.  No retained line contains a figure environment, an image-inclusion command or drawing code, so no image is delivered and the package declares no asset dependency.  Editorial formatting only, in addition: an AMS \texttt{amsart} preamble that restates the article's own declarations and macros used by the retained lines, namely the environments \texttt{theorem}, \texttt{lemma} on separate counters, together with the standard packages those lines need.  The retained environments print in this document as Theorem 1, Lemma 1 in order of appearance, whereas the article prints the same items with the numbering of its own full document.  The complete native source of the article as distributed in the arXiv e-print for this exact version is bundled unchanged as \texttt{sources/210759/source0.tex}.
\begin{theorem} \label{heisenberg}
   For  \(\mathfrak{h}(1)=\langle x,y,z \mid [x,y]=z\rangle \) on \(\mathbb{F}_{p}\) we have 
   $$\mathrm{sd}(\mathfrak{h}(1))=\frac{3p^3+12p^2+16p+16}{{(p^2+2p+4)^2}}.$$
\end{theorem}

\begin{equation} \label{permutewithA}
        \mathcal{Q}_{\mathcal{L}(L)}(A) =\{B \in \mathcal{L}  (L)  \mid  \ [A, \  B] \subseteq A+B \}. 
\end{equation}

\begin{lemma}\label{combinatorialformulat} Given $L$ finite dimensional Lie algebra on $\mathbb{F}_{p^n}$, we have
\[\label{sdviacharfun}
    \mathrm{sd}(L)=\frac{1}{|\mathcal{L}(L)|^2} \underset{A, B\in \mathcal{L}(L) }{\sum} \chi (A,B).
\]    
\end{lemma}

\begin{theorem} \label{sdforsl2}
 If $p$ is an odd prime and $q=p^n$, then $\mathrm{sd}(\mathfrak{sl}_2(\mathbb{F}_q))$ has the same type of growth of $\mathrm{sd}(\mathfrak{h}(1))$ in Theorem \ref{heisenberg} replacing the role of $q$ with that of $p$, that is,   $$\mathrm{sd}(\mathfrak{sl}_2(\mathbb{F}_q))=\frac{3q^3+12q^2+16q+16}{{(q^2+2q+4)^2}}.$$
\end{theorem}

\begin{proof}[Proof of Theorem \ref{sdforsl2}] 
We should remember that $q$ is a power of an odd prime $p$. Then, since 
 $\mathfrak{sl}_2(\mathbb{F}_q)=\Bigg \{ \begin{pmatrix}a&b \\ c&-a\end{pmatrix}\Big| a,b,c \in \mathbb{F}_q \Bigg \}$, it is easy to check that $ \Bigg \{e=\begin{pmatrix} 0 & 1 \\ 0 & 0 \end{pmatrix}, f=\begin{pmatrix} 0 & 0 \\ 1 & 0 \end{pmatrix}, h=\begin{pmatrix} 1 & 0 \\ 0 & -1 \end{pmatrix}\Bigg \}$ forms a basis for $\mathfrak{sl}_2(\mathbb{F}_q)$ with 
 \begin{equation} \label{bracket}
     [h, e] = 2e, \quad [h, f] = -2f , \quad [e, f] = h 
 \end{equation}
 and  $|\mathfrak{sl}_2(\mathbb{F}_q)|=q^3$.  Since every one dimensional subspace of $\mathfrak{sl}_2(\mathbb{F}_3)$ is a one dimensional Lie subalgebra, there are $\frac{q^3-1}{q-1}=q^2+q+1$ one dimensional Lie subalgebras such that
\begin{equation} \label{size of lines}
    q+1, \ \ \frac{q(q+1)}{2} \ \ \ \mbox{and}\ \ \ \frac{q(q-1)}{2}
\end{equation}
are  nilpotent, split semisimple  and  nonsplit semisimple,  respectively. Now,
using \eqref{bracket} we obtain that  $\langle e \rangle$, $\langle f \rangle$  and  $N_\mu=\langle h+\mu e -\mu^{-1}f\rangle$ for all $\mu \in\{1,2,\cdots, q-1\}$ are the one dimensional subalgebras coresponding to the nilpotent Lie subalgebras, and each  of them, along with the one dimensional subalgebras generated by split semisimple element, generates  two dimensional  subalgebras in $\mathfrak{sl}_2(\mathbb{F}_q)$. 
This means, if $S_i$ (with $ \ i \in\{1,2,\cdots, q-1\}$) is split semisimple, that
$$\mathfrak{b}_+=\langle \langle e \rangle, \langle h \rangle \rangle, \ \ \mathfrak{b}_-=\langle \langle f \rangle, \langle h \rangle \rangle, \ \
\mathfrak{b}_{\mu}= \langle N_\mu, S_i \rangle $$  are all the $q+1$ Borel subalgebras of $\mathfrak{sl}_2(\mathbb{F}_q)$, and one is nilpotent while $q$ split semisimple.  
This implies, including the trivial ideal and $\mathfrak{sl}_2(\mathbb{F}_q)$ itself, that
\begin{equation}\label{size of the lattice of sl2}
      |\mathcal{L}(\mathfrak{sl}_2(\mathbb{F}_q))|=2+(q^2+q+1)+(q+1)= q^2+2q+4 \ \ \mbox{and} \ \ |\mathcal{Q}_{\mathcal{L}(\mathfrak{sl}_2(\mathbb{F}_q))}(\mathcal{L}(\mathfrak{sl}_2(\mathbb{F}_q)))|= q+3.
\end{equation}
Now, we can apply \eqref{permutewithA} to obtain $|\mathcal{Q}_{\mathcal{L}(\mathfrak{sl}_2(\mathbb{F}_q))}(A)|$ for each $A \in \mathcal{L}(\mathfrak{sl}_2(\mathbb{F}_q))$.  So, we consider the following cases. 

\medskip
\textit{Case 1.} Assume $A \in \mathcal{Q}_{\mathcal{L}(\mathfrak{sl}_2(\mathbb{F}_q))}(\mathcal{L}(\mathfrak{sl}_2(\mathbb{F}_q)))$. 

\medskip
In this case it is obvious that $$\mathcal{Q}_{\mathcal{L}(\mathfrak{sl}_2(\mathbb{F}_q))}(A)=q^2+q+4.$$


\medskip
\textit{Case 2.} Let $A$ be  nilpotent subalgebra. 

\medskip
Since each nilpotent Lie algebra in $\mathfrak{sl}_2(\mathbb{F}_q)$ along with split semisimple are enough to generate a Borel subalgebra,  each of them can not be contained in two different Borel subalgebras. In fact,  the intersection of  two Borel  subalgebras is always a  one dimensional subalgebra which is generated by a semisimple element. This implies that each $q+1$ nilpotent Lie algebra of dimension one permutes with exactly $q+1$ one dimensional subalgebras contained in the same Borel subalgebra. In other words, \eqref{permutewithA}  implies 
$$\mathcal{Q}_{\mathcal{L}(\mathfrak{sl}_2(\mathbb{F}_q))}(A)=(q+1)+(q+3)=2q+4.$$ 

\medskip
\textit{Case 3.} Let $A$ be split semisimple subalgebra.

\medskip
In this case, each $A$ is the intersection of two different Borel subalgebras.  This means  $A$ permutes with $q+1$ subalgebras in the first Borel subalgebra and with the remaining $q$ subalgebras contained in the second Borel subalgebra. Thus,
$$\mathcal{Q}_{\mathcal{L}(\mathfrak{sl}_2(\mathbb{F}_q))}(A)=(q+1)+q=(2q+1)+(q+3)=3q+4.$$

\medskip
\textit{Case 4.} Let $A$ be  nonsplit semisimple subalgebra.

\medskip
Since there are exactly $\frac{q(q-1)}{2}$ maximal one dimensional subalgebras in $\mathfrak{sl}_2(\mathbb{F}_q)$, which do not generate a Borel subalbegra, then each of them do not permute with every other one dimensional subalgebra. This implies $$\mathcal{Q}_{\mathcal{L}(\mathfrak{sl}_2(\mathbb{F}_q))}(A)=q+4.$$ 
Therefore, by applying Lemma \ref{combinatorialformulat} with \eqref{size of lines} and \eqref{size of the lattice of sl2} we obtain
$$(q^2+2q+4)^2\mathrm{sd}(\mathfrak{sl}_2(\mathbb{F}_q)))=(q+1)(q^2+2q+4)+(q+3)(2q+4)+\Big(\frac{q(q+1)}{2}\Big)(3q+4)+\Big(\frac{q(q-1)}{2}\Big)(q+4).$$
This completes the proof.
\end{proof}

\end{document}
